% language=us

\environment context-2022-style

\startdocument
  [title={CTX Math},
   banner={the updated engine},
   location={context\enspace {\bf 2022}\enspace meeting}]

\starttitle[title=A little history]

\startitemize
    \startitem
        When we started with \LUATEX\ there was little known about
        \UNICODE\ and \OPENTYPE\ math. References were \MSWORD\ and
        Cambria.
    \stopitem
    \startitem
        So we ended up with a hybrid math machinery that has to
        support traditional and new math fonts.
    \stopitem
    \startitem
        Later the mixed code paths were separated a bit more and some operating
        modes were introduced. The biggest issue is in the way traditional \TEX\
        fonts (ab)use italic correction and some of that (inconsistently) made it
        into modern fonts.
    \stopitem
    \startitem
        Once fonts started showing up it became clear that we could
        never really do it well. In \CONTEXT\ we added support for
        font specific runtime patching.
    \stopitem
    \startitem
        In \LUAMETATEX\ we were free from existing assumptions and didn't need
        to take other \quote {frozen in time assumptions} into account.
    \stopitem
    \startitem
        So, after a lot of extra control was added in the engine, Mikael and I set
        off to use it and add more.
    \stopitem
    \startitem
        Once we were satisfied about the mix of engine features and runtime
        font patching some control was removed from the engine.
    \stopitem
    \startitem
        And while we were getting fonts under control we were also exploring
        fundamental extension to the rendering (like spacing and positioning).
    \stopitem
    \startitem
        This also permitted us to review the \CONTEXT\ math mechanisms, simplify some
        of the implementation and extend them. Optional code paths were removed.
    \stopitem
    \startitem
        So where are we after a few man|-|years of work?
    \stopitem
\stopitemize

\stoptitle

\starttitle[title=The process]

\startitemize
    \startitem
        The math engine is mostly separated from the rest: the main function is
        \type {mlist_to_hlist} which is also called recursively when a formula gets
        done.
    \stopitem
    \startitem
        We have yet unprocessed nodes, called noads: atoms, fractions, radicals,
        fences and accents. Plus some style and choice related ones.
    \stopitem
    \startitem
        All become regular glyph, list, glue, penalty nodes.
    \stopitem
    \startitem
        We support boundaries and discretionaries in math and math in
        discretionaries.
    \stopitem
    \startitem
        Rendering, spacing, penalties, unrolling, etc.\ are all controlled
        by classes and their related options.
    \stopitem
    \startitem
        There are new concepts, like dictionaries, direct fences, generic
        radicals, vertical stacking (for arrows), etc.
    \stopitem
    \startitem
        There is also plenty of (old and upgraded) handling at the \CONTEXT\
        end.
    \stopitem
    \startitem
        Mikaels presentation gives a nice overview.
    \stopitem
\stopitemize

\stoptitle

\starttitle[title=Maybe]

\startitemize
    \startitem
        Remove display math because it is basically inline math (in display
        style) with some spacing around and a formula number attached.
    \stopitem
    \startitem
        Remove old font related character definition primitives because we can
        assume \OPENTYPE\ fonts to be used (once we have a proper Euler).
    \stopitem
    \startitem
        Get rid of the \type {U} prefix.
    \stopitem
    \startitem
        Use a generic parameter setter instead of the plenty ones (using symbolic
        names for indices).
    \stopitem
\stopitemize

{\em See Mikaels various math related presentations at various meetings about
what we did. The \quote {Math in \CONTEXT} manual gives a complete overview.}

\stoptitle

% Mikaels talk and demos.
% Our math-unicode document.

\stopdocument
